% Intended LaTeX compiler: pdflatex
\documentclass[a4paper,12pt]{article}
\usepackage[utf8]{inputenc}
\usepackage[T1]{fontenc}
\usepackage{graphicx}
\usepackage{longtable}
\usepackage{wrapfig}
\usepackage{rotating}
\usepackage[normalem]{ulem}
\usepackage{amsmath}
\usepackage{amssymb}
\usepackage{capt-of}
\usepackage{hyperref}
\usepackage{\string~/.emacs.d/common}
\author{mahmood}
\date{\today}
\title{maxima in lisp}
\hypersetup{
 pdfauthor={mahmood},
 pdftitle={maxima in lisp},
 pdfkeywords={},
 pdfsubject={},
 pdfcreator={Emacs 28.2 (Org mode 9.5.5)}, 
 pdflang={English}}
\begin{document}

\maketitle
\begin{note}
this document and others can be found on my website\\
\url{https://mahmoodsheikh36.github.io/}
\end{note}
\begin{note}
i ended up giving up on using maxima as a visualization/symbolic math computation library for elisp, it wasnt made with this usecase in mind i guess, and so its an utter nightmare to lisp with it\\
\end{note}
\href{20230510181855-maxima.org}{maxima} is written in \href{20230215220838-lisp.org}{lisp} and so it can be used as a modified/extended lisp runtime\\
\section{how to}
\label{sec:org062e185}
to start the lisp runtime use the command:\\
\lstset{language=bash,label= ,caption= ,captionpos=b,numbers=none}
\begin{lstlisting}
rmaxima -r 'to_lisp();'
\end{lstlisting}
can also run this in \href{emacs.org}{emacs}' slime/sly with \texttt{(sly "rmaxima -r to\_lisp();")}\\
alternatively, we can clone maxima's source code to the subdirectory \texttt{local-projects} of the quicklisp installation directory (usually \texttt{\textasciitilde{}/quicklisp} unless you've modified it), follow the instructions at \url{https://sourceforge.net/p/maxima/code/ci/master/tree/INSTALL.lisp\#l66} to compile maxima and you'll be able to load it as a library:\\
\lstset{language=Lisp,label= ,caption= ,captionpos=b,numbers=none}
\begin{lstlisting}
(ql:quickload "maxima")
\end{lstlisting}
but then you gotta prefix functions with \texttt{maxima::}, e.g.\\
\lstset{language=Lisp,label= ,caption= ,captionpos=b,numbers=none}
\begin{lstlisting}
(maxima::displa (maxima::$integrate #$2/(3*x^5)$ 'x))
\end{lstlisting}

\begin{verbatim}
2 x
----
   5
3 x
\end{verbatim}


unless you enter the library itself then you dont need the prefix:\\
\lstset{language=Lisp,label= ,caption= ,captionpos=b,numbers=none}
\begin{lstlisting}
(in-package :maxima)
\end{lstlisting}
loading maxima as a library works most of the time, but it causes some problems, for example we cant use \texttt{draw} when maxima is loaded as a library, thats why using \texttt{rmaxima} is better\\
\section{misc}
\label{sec:orgac397ad}
to write an expression in infix syntax we do:\\
\lstset{language=Lisp,label= ,caption= ,captionpos=b,numbers=none}
\begin{lstlisting}
(print #$10/13$)
\end{lstlisting}

\begin{verbatim}

((MAXIMA::RAT MAXIMA::SIMP) 10 13) 
\end{verbatim}


to display math using ascii art we use the function \texttt{displa} (short for \textbf{display})\\
\lstset{language=Lisp,label= ,caption= ,captionpos=b,numbers=none}
\begin{lstlisting}
(maxima::displa #$10/13$)
\end{lstlisting}

\begin{verbatim}
10
--
13
\end{verbatim}
\section{maxima expression to lisp expression}
\label{sec:orgcb5abc7}
this function helps turn maxima expressions into their corresponding lisp expression\\
\lstset{language=Lisp,label= ,caption= ,captionpos=b,numbers=none}
\begin{lstlisting}
(defun lisp-form (macsyma-string)
  (maxima::macsyma-read-string (concatenate 'string macsyma-string "$")))
\end{lstlisting}

example:\\
\lstset{language=Lisp,label= ,caption= ,captionpos=b,numbers=none}
\begin{lstlisting}
(print (lisp-form "diff(sin(x),x)"))
\end{lstlisting}

\begin{verbatim}

((MAXIMA::$DIFF) ((MAXIMA::%SIN) MAXIMA::$X) MAXIMA::$X) 
\end{verbatim}


some expressions will return symbols instead of functions so they cannot be run directly, e.g.\\
\lstset{language=Lisp,label= ,caption= ,captionpos=b,numbers=none}
\begin{lstlisting}
(print (lisp-form "sin(10d0)"))
\end{lstlisting}

\begin{verbatim}

((MAXIMA::%SIN) 10.0) 
\end{verbatim}


so to evaluate these expressions we can use \texttt{meval*}:\\
\lstset{language=Lisp,label= ,caption= ,captionpos=b,numbers=none}
\begin{lstlisting}
(print (lisp-form "sin(10d0)"))
(print (maxima::meval* (lisp-form "sin(10d0)")))
(print (maxima::meval* '((maxima::%sin) 10.0)))
\end{lstlisting}

\begin{verbatim}

((MAXIMA::%SIN) 10.0) 
-0.5440211108893698 
-0.5440211108893698 
\end{verbatim}


although functions like \texttt{\$sin} do exist and can be used instead of symbols like \texttt{\%sin}\\
taken from \url{https://github.com/livelisp/livewlisp/blob/main/maxima-tutorial.txt}\\
\section{integration}
\label{sec:org70cb455}
the main function for integration is \texttt{\$integrate} (or \texttt{sinint})\\
\lstset{language=Lisp,label= ,caption= ,captionpos=b,numbers=none}
\begin{lstlisting}
(maxima::displa (maxima::$integrate #$2/(3*x^2)$ 'x))
\end{lstlisting}

\begin{verbatim}
2 x
----
   2
3 x
\end{verbatim}


\lstset{language=Lisp,label= ,caption= ,captionpos=b,numbers=none}
\begin{lstlisting}
(maxima::displa (maxima::$integrate '((maxima::%cos) x) 'x))
\end{lstlisting}

\begin{verbatim}
sin(x)
\end{verbatim}

\section{lists}
\label{sec:orgadeb7e9}
i havent a better approach yet\\
\lstset{language=Lisp,label= ,caption= ,captionpos=b,numbers=none}
\begin{lstlisting}
(maxima::displa (maxima::meval* `((maxima::mlist) 2 3 5)))
\end{lstlisting}

lisp list to maxima list:\\
\lstset{language=Lisp,label= ,caption= ,captionpos=b,numbers=none}
\begin{lstlisting}
(defun list->mlist (list)
  (let ((mlist (maxima::meval* `((maxima::mlist)))))
    (loop for expr in (reverse list) ;; its reversed in maxima dunno why
          do (setf
              mlist
              (maxima::meval* `((maxima::$append) ((maxima::mlist) ,expr) ,mlist))))
    mlist))
\end{lstlisting}
example:\\
\lstset{language=Lisp,label= ,caption= ,captionpos=b,numbers=none}
\begin{lstlisting}
(print (list->mlist '(1 2 3)))
(maxima::displa (list->mlist '(1 2 3)))
\end{lstlisting}

\section{plotting}
\label{sec:org91fff8f}
we need to initialize some variables (like \texttt{*maxima-tempdir*}) to be able to plot, this can be done automatically using \texttt{initialize-runtime-globals}\\
\lstset{language=Lisp,label= ,caption= ,captionpos=b,numbers=none}
\begin{lstlisting}
(maxima::initialize-runtime-globals)
\end{lstlisting}
to plot a set of points (discrete plot) using \texttt{gnuplot}\\
\lstset{language=Lisp,label= ,caption= ,captionpos=b,numbers=none}
\begin{lstlisting}
(maxima::$plot2d
 '((maxima::mlist)
   maxima::$discrete
   ((maxima::mlist) 1 2 3) ((maxima::mlist) 1 2 3)))
\end{lstlisting}
this can be (somewhat) simplified using our function \texttt{list->mlist} (see above)\\
\lstset{language=Lisp,label= ,caption= ,captionpos=b,numbers=none}
\begin{lstlisting}
(maxima::$plot2d
 `((maxima::mlist)
   ((maxima::mlist) maxima::$discrete
                    ,(list->mlist '(1 2 3)) ,(list->mlist '(1 2 3)))
   ((maxima::mlist) maxima::$discrete
                    ,(list->mlist '(1 2 3)) ,(list->mlist '(1 5 3)))))
\end{lstlisting}
we can draw multiple plots (this only works when running lisp using \texttt{maxima}, see the \textbf{how to} section):\\
\lstset{language=Lisp,label= ,caption= ,captionpos=b,numbers=none}
\begin{lstlisting}
(let ((scene1 '((MAXIMA::$GR2D)
                ((MAXIMA::MEQUAL) MAXIMA::$TITLE "first plot")
                ((MAXIMA::MEQUAL) MAXIMA::$NTICKS 300)
                ((MAXIMA::$PARAMETRIC)
                 ((MAXIMA::MTIMES) 2 ((MAXIMA::%COS) MAXIMA::$T))
                 ((MAXIMA::MTIMES) 5 ((MAXIMA::%SIN) MAXIMA::$T)) MAXIMA::$T 0
                 ((MAXIMA::MTIMES) 2 MAXIMA::$%PI))))
      (scene2 `((MAXIMA::$GR2D)
                ((MAXIMA::MEQUAL) MAXIMA::$TITLE "second plot")
                ((MAXIMA::MEQUAL) MAXIMA::$NTICKS 300)
                ((maxima::$points) ((mlist) 1 2 3) ((mlist) 1 2 3)))))
  (maxima::meval* `((MAXIMA::$DRAW) ,scene1 ,scene2 ((MAXIMA::MEQUAL) MAXIMA::$COLUMNS 2))))
\end{lstlisting}
we implement some abstraction over this to make it less explicit:\\
\lstset{language=Lisp,label= ,caption= ,captionpos=b,numbers=none}
\begin{lstlisting}
(defun plot-points (x-values y-values)
  (maxima::$plot2d
   `((maxima::mlist)
     maxima::$discrete
     ,(list->mlist x-values) ,(list->mlist y-values))))
\end{lstlisting}
\end{document}